% 第 16 章：子类型化的元理论
% Source: source/split/chapters/ch16-metatheory-of-subtyping_p200-p211.pdf
% Original print pages: 209--220
% Status: 人工总审中

\chapter{子类型化的元理论}
\label{chap:16}

\section*{概述}
\phantomsection
\label{sec:16-overview}

上一章定义的带子类型化简单类型 lambda 演算还不适合直接实现。与此前的演算
不同，这套规则不是\termfirst{语法导向的}{syntax directed}：给定一个待检查的
结论时，无法总是唯一确定应采用哪条规则，也无法总是由结论直接确定其前提。
因此，不能直接把这些推导规则倒着执行，将其作为类型检查算法。主要问题来自
类型关系中的包容规则 \textsc{T-Sub} 和子类型关系中的传递规则 \textsc{S-Trans}。

\textsc{T-Sub} 的结论只把项写成不带任何结构限制的元变量 \(t\)：
\[
\inferrule*[right=\textsc{T-Sub}]
 {\Gamma\vdash t:S\\S\subtype T}
 {\Gamma\vdash t:T}
\]
其他类型规则都要求项具有某种具体形式：\textsc{T-Abs} 只适用于 lambda
抽象，\textsc{T-Var} 只适用于变量，等等。\textsc{T-Sub} 却能用于任何项。
因此，给定一个待计算类型的项 \(t\) 时，既可能选择 \textsc{T-Sub}，也可能
选择结论与 \(t\) 的句法形式相符的那条规则。

\textsc{S-Trans} 也有同样的问题：它的结论与所有其他子类型规则的结论重叠。
\[
\inferrule*[right=\textsc{S-Trans}]
 {S\subtype U\\U\subtype T}
 {S\subtype T}
\]
\(S\) 与 \(T\) 都是不带结构限制的元变量，所以任何子类型判断都可能以
\textsc{S-Trans} 收尾。若把规则机械地从结论向前提执行，算法无法知道应该
尝试传递规则，还是应该尝试另一条与输入类型结构相符的规则。

\textsc{S-Trans} 还有第二个问题：两个前提都出现了结论中没有的 \(U\)。
从结论反推前提时，必须先猜一个 \(U\)，再检查 \(S\subtype U\) 和
\(U\subtype T\)。可供猜测的类型有无穷多个，这种搜索方式无法成为实用算法。

\textsc{S-Refl} 的结论也与其他子类型规则重叠，不过问题较轻：它没有前提；
输入两端相同时可以立即成功。即便如此，它仍使规则不是语法导向的。

解决方法是用两套语法导向的关系替换普通的、也称
\termfirst{声明式}{declarative}的子类型关系和类型关系：新关系分别称为
\termfirst{算法子类型关系}{algorithmic subtyping relation}和
\termfirst{算法类型关系}{algorithmic typing relation}。随后证明两种表述
实际上等价：算法规则可导出的子类型判断恰好也是声明式规则可导出的判断；
一个项能由算法规则赋型，当且仅当它能由声明式规则赋型。

\taplsecref{16.1}{sec:16.1}构造算法子类型关系，
\taplsecref{16.2}{sec:16.2}构造算法类型关系。
\taplsecref{16.3}{sec:16.3}加入\cref{fig:8-1}中的布尔值与条件式。为条件式
赋型时，还需要计算两个分支类型的最小公共超类型；这个类型称为二者的并。
\taplsecref{16.4}{sec:16.4}讨论最小类型
\(\tapltype{Bot}\)。\footnote{本章研究\cref{fig:15-1,fig:15-3}中的演算。
原书相应 OCaml 实现为 \texttt{rcdsub}；\taplsecref{16.3}{sec:16.3}
加入布尔值与条件式，相应实现为 \texttt{joinsub}；
\taplsecref{16.4}{sec:16.4}加入 \(\tapltype{Bot}\)，相应实现为
\texttt{bot}。}

\section{算法子类型关系}
\label{sec:16.1}

带子类型化语言的实现必须能够检查一个类型是否为另一个类型的子类型。例如，
类型检查器遇到应用 \(t_1\,t_2\)，并算出
\(t_1:T\to U\)、\(t_2:S\) 时，需要判断 \(S\subtype T\) 能否由
\cref{fig:15-1,fig:15-3}中的规则导出。

子类型检查器实际判断的是另一个关系
\[
  \algturnstile S\subtype T
\]
读作“\(S\) 在算法关系下是 \(T\) 的子类型”。算法可以根据 \(S\) 与 \(T\)
的类型结构选择规则，再递归检查它们的组成部分。虽然算法关系使用的规则与
声明式关系不同，但二者给出的判断结果完全相同：声明式关系能够证明
\(S\subtype T\)，当且仅当算法关系也能判定 \(S\subtype T\)。两者最明显的
区别是，算法关系没有 \textsc{S-Trans} 和 \textsc{S-Refl}。

不过，在删去传递规则之前，先要重组记录子类型规则。如
\taplsecref{15.2}{sec:15.2}所见，宽度、深度和
排列推导必须借助传递性拼接。将三条规则合并为一条宏规则：
\[
\inferrule*[right=\textsc{S-Rcd}]
 {\{l_i\mid i\in1..n\}\subseteq\{k_j\mid j\in1..m\}\\
  k_j=l_i\Longrightarrow S_j\subtype T_i}
 {\{k_j:S_j\}^{j\in1..m}\subtype
  \{l_i:T_i\}^{i\in1..n}}
\]
第二个前提表示：对目标类型中的每个标签 \(l_i\)，都在源类型中找到同名字段
\(k_j\)，并检查对应字段类型 \(S_j\subtype T_i\)。

\begin{lemma}
\label{lem:16.1.1}
若 \(S\subtype T\) 能由包含 \textsc{S-RcdDepth}、
\textsc{S-RcdWidth} 和 \textsc{S-RcdPerm}、但不含 \textsc{S-Rcd} 的
规则导出，那么也能由使用 \textsc{S-Rcd}、但不使用前三条规则的系统导出；
反之亦然。
\end{lemma}
\begin{proof}
对推导作直接归纳。
\end{proof}

\cref{lem:16.1.1}允许用 \textsc{S-Rcd} 取代三条分散的记录规则。
\cref{fig:16-1}汇总重组后的系统。

\begin{figure}[!ht]
\centering
\begin{minipage}{0.98\linewidth}
\small
\taplfigureheader{{\Large \(<:\)}}
  {\cref{fig:15-1,fig:15-3}}
\[
\inferrule*[right=\textsc{S-Refl}]{ }{S\subtype S}
\]
\[
\inferrule*[right=\textsc{S-Trans}]
 {S\subtype U\\U\subtype T}{S\subtype T}
\]
\[
\inferrule*[right=\textsc{S-Top}]{ }{S\subtype\tapltype{Top}}
\]
\[
\inferrule*[right=\textsc{S-Arrow}]
 {T_1\subtype S_1\\S_2\subtype T_2}
 {S_1\to S_2\subtype T_1\to T_2}
\]
\[
\inferrule*[right=\textsc{S-Rcd}]
 {\{l_i\mid i\in1..n\}\subseteq\{k_j\mid j\in1..m\}\\
  k_j=l_i\Longrightarrow S_j\subtype T_i}
 {\{k_j:S_j\}^{j\in1..m}\subtype
  \{l_i:T_i\}^{i\in1..n}}
\]
\end{minipage}
\caption{带记录的子类型关系（紧凑版本）}
\label{fig:16-1}
\end{figure}
\FloatBarrier

下面证明，在\cref{fig:16-1}的系统中，自反规则和传递规则并非必需。

\begin{lemma}
\label{lem:16.1.2}
\begin{enumerate}
\item 对每个类型 \(S\)，不用 \textsc{S-Refl} 也能导出
  \(S\subtype S\)。
\item 若判断 \(S\subtype T\) 能由\cref{fig:16-1}中的规则导出，那么还存在
  一份不使用 \textsc{S-Trans} 的推导，其结论仍为 \(S\subtype T\)。
\end{enumerate}
\end{lemma}
\begin{proof}
练习（推荐，\(\star\star\star\)）。
\end{proof}
\taplauthorsolutionlink{sol:author:16.1.2}{16.1.2}

\begin{exercise}[\(\star\)]
\label{ex:16.1.3}
若加入类型 \(\tapltype{Bool}\)，这两项性质需要怎样修改？
\end{exercise}
\taplauthorsolutionlink{sol:author:16.1.3}{16.1.3}

\begin{definition}
\label{def:16.1.4}
算法子类型关系是由
\cref{fig:16-2}中的规则闭合生成的最小类型关系。
\end{definition}

\begin{figure}[!ht]
\centering
\begin{minipage}{0.98\linewidth}
\small
\taplfigureheaderplain{\(\algturnstile S\subtype T\)}
  {\tapltranslateditalic{算法子类型关系}}
\[
\inferrule*[right=\textsc{SA-Top}]
 { }{\algturnstile S\subtype\tapltype{Top}}
\]
\[
\inferrule*[right=\textsc{SA-Rcd}]
 {\{l_i\mid i\in1..n\}\subseteq\{k_j\mid j\in1..m\}\\
  k_j=l_i\Longrightarrow\algturnstile S_j\subtype T_i}
 {\algturnstile\{k_j:S_j\}^{j\in1..m}\subtype
  \{l_i:T_i\}^{i\in1..n}}
\]
\[
\inferrule*[right=\textsc{SA-Arrow}]
 {\algturnstile T_1\subtype S_1\\
  \algturnstile S_2\subtype T_2}
 {\algturnstile S_1\to S_2\subtype T_1\to T_2}
\]
\end{minipage}
\caption{算法子类型关系}
\label{fig:16-2}
\end{figure}
\FloatBarrier

算法规则的\emph{可靠性}是指：它导出的每个判断也能由声明式规则导出；换言之，
算法不会错误接受声明式规则无法证明的判断。算法规则的\emph{完备性}是指：
声明式规则能够导出的每个判断，算法规则也能导出；换言之，声明式系统认可的
判断不会被算法遗漏。

\begin{proposition}[可靠性与完备性]
\label{prop:16.1.5}
\[
  S\subtype T
  \quad\Longleftrightarrow\quad
  \algturnstile S\subtype T
\]
\end{proposition}
\begin{proof}
两个方向都对推导作归纳，分别使用\cref{lem:16.1.1,lem:16.1.2}。
\end{proof}

现在可以把这些算法规则逐条转写为子类型检查算法：
\begin{lstlisting}
subtype(S, T) =
  if T = Top then true
  else if S = S1 -> S2 and T = T1 -> T2 then
    subtype(T1, S1) and subtype(S2, T2)
  else if S = {kj:Sj} and T = {li:Ti} then
    every target label li occurs as some kj
    and subtype(Sj, Ti) for every matching field
  else false
\end{lstlisting}
具体的 OCaml 实现见\chapref{17}。

最后还要证明算法对每一对输入都在有限时间内返回 \texttt{true} 或
\texttt{false}，即它是全函数。

\begin{proposition}[终止性]
\label{prop:16.1.6}
若 \(\algturnstile S\subtype T\) 可导出，则
\(\mathsf{subtype}(S,T)\) 返回 \texttt{true}；否则返回 \texttt{false}。
\end{proposition}
\begin{proof}
首先，若算法子类型判断 \(\algturnstile S\subtype T\) 可导出，则对这份
推导作归纳，可以验证 \(\mathsf{subtype}(S,T)\) 会沿着相应的算法分支执行并
返回 \texttt{true}。反过来，若 \(\mathsf{subtype}(S,T)\) 返回
\texttt{true}，根据它所经过的分支递归构造推导，也能得到
\(\algturnstile S\subtype T\)。

现在考虑该判断不可导出的情况。根据上面的结论，算法不可能返回
\texttt{true}；要证明它一定返回 \texttt{false}，还需排除无限递归。把类型
的大小理解为其语法树中构造节点的数量。算法每次递归调用时，传入的都是原类型的
组成部分，因此两个输入类型的大小之和都会严格减小。这个和始终是正整数，所以
不可能存在无限的递归调用序列。算法必定终止；既然不能返回 \texttt{true}，
就只能返回 \texttt{false}。
\end{proof}

\cref{prop:16.1.5,prop:16.1.6}共同说明
\(\mathsf{subtype}\) 是声明式子类型关系的
\termfirst{判定过程}{decision procedure}。

也可以一开始就把算法关系当作子类型关系的正式定义，完全不提声明式版本。
不过，这并不能省去多少工作：为了证明依赖子类型关系的类型系统具有良好性质，
仍需证明算法子类型关系具有自反性和传递性，工作量与本节相近。实际语言定义
有时确实直接采用算法式类型规则；\chapref{19}就是一个例子。

\section{算法类型关系}
\label{sec:16.2}

处理好子类型关系后，还要以同样方式处理类型关系。如
\hyperref[sec:16-overview]{本章概述}所见，唯一不语法导向的类型规则是
\textsc{T-Sub}。它不能简单删掉；必须先找出包容规则在类型推导中不可替代的
用途，再把这种用途嵌入其他类型规则。

在函数应用中，包容规则负责连接函数期望的参数类型与实参的实际类型。例如
\[
  (\lambda r:\{x:\tapltype{Nat}\}.\,r.x)\,
  \{x=0,y=1\}
\]
离开包容规则便无法赋型。

稍微出人意料的是，这也是包容规则唯一不可替代的用途。在其他位置，可以重排
类型推导，把 \textsc{T-Sub} 的使用逐步推迟到更靠近最终结论的位置。

\translatornote{这里重排 \textsc{T-Sub}，是为了从声明式类型规则中提取可直接
执行的类型检查算法。\textsc{T-Sub} 可以在推导的任意位置把已经得到的类型提升为
任意超类型；如果它散落在推导各处，类型检查器就无法只根据项的语法确定下一步。
把它尽量移到推导末端后，末端的 \textsc{T-Sub} 只会把已经算出的类型换成信息
更少的超类型，算法可以省略这一步，保留更小、更精确的类型。唯一无法完全消除的
情形是函数应用中实参类型与参数类型的匹配；这项必要的子类型检查随后会被吸收到
增强后的应用规则 \textsc{TA-App} 中。下面的推导重排正是在说明，经过这样的整理，
原来能够赋型的项仍然能够通过类型检查。}

先看以 \textsc{T-Abs} 结束、而其直接子推导以 \textsc{T-Sub} 结束的推导：
\[
\inferrule*[right=\textsc{T-Abs}]
 {\inferrule*[right=\textsc{T-Sub}]
   {\Gamma,x:S_1\vdash s_2:S_2\\S_2\subtype T_2}
   {\Gamma,x:S_1\vdash s_2:T_2}}
 {\Gamma\vdash\lambda x:S_1.\,s_2:S_1\to T_2}
\]
可以把包容规则移到抽象规则之后：
\[
\inferrule*[right=\textsc{T-Sub}]
 {\inferrule*[right=\textsc{T-Abs}]
  {\Gamma,x:S_1\vdash s_2:S_2}
   {\Gamma\vdash\lambda x:S_1.\,s_2:S_1\to S_2}
  \\
  \inferrule*[right=\textsc{S-Arrow}]
   {\inferrule*[right=\textsc{S-Refl}]
     { }{S_1\subtype S_1}
    \\
    S_2\subtype T_2}
   {S_1\to S_2\subtype S_1\to T_2}}
 {\Gamma\vdash\lambda x:S_1.\,s_2:S_1\to T_2}
\]

应用规则有两个直接子推导，其中任意一个都可能以 \textsc{T-Sub} 结束。先看
包容规则出现在左侧函数子推导末尾的情形。重排前的推导为：
\[
\inferrule*[right=\textsc{T-App}]
 {\inferrule*[right=\textsc{T-Sub}]
   {\Gamma\vdash s_1:S_{11}\to S_{12}
    \\
    \inferrule*[right=\textsc{S-Arrow}]
     {T_{11}\subtype S_{11}\\S_{12}\subtype T_{12}}
     {S_{11}\to S_{12}\subtype T_{11}\to T_{12}}}
   {\Gamma\vdash s_1:T_{11}\to T_{12}}
  \\
  \Gamma\vdash s_2:T_{11}}
 {\Gamma\vdash s_1\,s_2:T_{12}}
\]
这里先用 \textsc{T-Sub} 把函数类型从 \(S_{11}\to S_{12}\) 提升为
\(T_{11}\to T_{12}\)，再应用于 \(T_{11}\) 型实参。子类型反演给出
\(T_{11}\subtype S_{11}\) 和 \(S_{12}\subtype T_{12}\)。重排后，前一项
用来把实参类型提升到 \(S_{11}\)，后一项则在应用完成后把结果从
\(S_{12}\) 提升到 \(T_{12}\)：
\[
\inferrule*[right=\textsc{T-Sub}]
 {\inferrule*[right=\textsc{T-App}]
   {\Gamma\vdash s_1:S_{11}\to S_{12}
    \\
    \inferrule*[right=\textsc{T-Sub}]
     {\Gamma\vdash s_2:T_{11}\\T_{11}\subtype S_{11}}
     {\Gamma\vdash s_2:S_{11}}}
   {\Gamma\vdash s_1\,s_2:S_{12}}
  \\
  S_{12}\subtype T_{12}}
 {\Gamma\vdash s_1\,s_2:T_{12}}
\]
这里发生了两种移动。原来 \textsc{S-Arrow} 右侧的前提
\(S_{12}\subtype T_{12}\) 被下移到整个应用推导的末端，用来提升最终结果
类型；原来 \textsc{T-App} 右侧、用于推导实参
\(\Gamma\vdash s_2:T_{11}\) 的子推导，则被嵌入新的实参推导中，再由
\(T_{11}\subtype S_{11}\) 把实参提升为函数真正接受的类型。

再看包容规则出现在 \textsc{T-App} 的右侧实参子推导末尾的情形。重排前的
推导为：
\[
\inferrule*[right=\textsc{T-App}]
 {\Gamma\vdash s_1:T_{11}\to T_{12}
  \\
  \inferrule*[right=\textsc{T-Sub}]
   {\Gamma\vdash s_2:T_2\\T_2\subtype T_{11}}
   {\Gamma\vdash s_2:T_{11}}}
 {\Gamma\vdash s_1\,s_2:T_{12}}
\]
若仍要移动右侧实参子推导末尾的这次 \textsc{T-Sub}，唯一可行的做法是把它
转移到左侧函数子推导中。它与前一种重排的方向恰好相反：前面把“使函数
与实参的类型相匹配”这项调整从函数一侧移到了实参一侧，这里则将它移回
函数一侧。具体做法是利用参数类型的逆变性，把函数的参数类型从 \(T_{11}\)
改为实参原有的 \(T_2\)：
{\small
\[
\inferrule*[right=\textsc{T-App}]
 {\inferrule*[right=\textsc{T-Sub}]
   {\Gamma\vdash s_1:T_{11}\to T_{12}
    \\
    \inferrule*[right=\textsc{S-Arrow}]
     {T_2\subtype T_{11}
      \\
      \inferrule*[right=\textsc{S-Refl}]
       { }{T_{12}\subtype T_{12}}}
     {T_{11}\to T_{12}\subtype T_2\to T_{12}}}
   {\Gamma\vdash s_1:T_2\to T_{12}}
  \quad
  \Gamma\vdash s_2:T_2}
 {\Gamma\vdash s_1\,s_2:T_{12}}
\]
}
由此可见，这次包容可以在应用的两个前提之间移动，却无法彻底消除：既可以把
实参提升到函数要求的参数类型，也可以把函数调整为接受实参的实际类型。为了得到
统一的标准形式，下面固定采用“提升实参类型”的写法，把这次包容保留在应用的
右侧实参子推导末尾。

若一条 \textsc{T-Sub} 的直接子推导也以 \textsc{T-Sub} 结束，则两次使用
可以合并。重排前的推导为：
\[
\inferrule*[right=\textsc{T-Sub}]
 {\inferrule*[right=\textsc{T-Sub}]
   {\Gamma\vdash s:S\\S\subtype U}
   {\Gamma\vdash s:U}
  \\
  U\subtype T}
 {\Gamma\vdash s:T}
\]
由 \(S\subtype U\) 和 \(U\subtype T\) 先以 \textsc{S-Trans} 得到
\(S\subtype T\)，便可把相邻的两次 \textsc{T-Sub} 合并为一次：
\[
\inferrule*[right=\textsc{T-Sub}]
 {\Gamma\vdash s:S
  \\
  \inferrule*[right=\textsc{S-Trans}]
   {S\subtype U\\U\subtype T}
   {S\subtype T}}
 {\Gamma\vdash s:T}
\]

\begin{exercise}[\(\star\star\), \(\nrightarrow\)]
\label{ex:16.2.1}
完成上述实验：说明当 \textsc{T-Sub} 出现在 \textsc{T-Rcd} 或
\textsc{T-Proj} 之前时，怎样作类似重排。
\end{exercise}
\taplsolutionlink{sol:translator:16.2.1}{16.2.1}

固定上述重排方向后，反复应用这些变换，可以把任意类型推导整理成一种特殊形式：
\textsc{T-Sub} 只会出现在应用的右侧实参子推导末尾，以及整个推导的最后一步。
删掉最末端那次包容仍会为同一项留下一个类型，只是这个类型可能更小，也就是
携带更多信息。现在只有应用中的包容无法消除。把它吸收到应用规则中：
\[
\inferrule
 {\Gamma\vdash t_1:T_{11}\to T_{12}\\
  \Gamma\vdash t_2:T_2\\
  T_2\subtype T_{11}}
 {\Gamma\vdash t_1\,t_2:T_{12}}
\]
便能替换每个“应用之前先包容”的子推导，并完全删去 \textsc{T-Sub}。这条
增强后的应用规则仍然语法导向，因为结论中的项明确是一项应用。

\begin{definition}
\label{def:16.2.2}
算法类型关系是由
\cref{fig:16-3}中的规则闭合生成的最小关系。写作
\(\Gamma\algturnstile t:T\)，以区别于声明式判断 \(\Gamma\vdash t:T\)。
\end{definition}

\begin{figure}[!ht]
\centering
\begin{minipage}{0.98\linewidth}
\small
\taplfigureheaderplain{\(\Gamma\algturnstile t:T\)}
  {\tapltranslateditalic{算法类型关系}}
\[
\inferrule*[right=\textsc{TA-Var}]
 {x:T\in\Gamma}{\Gamma\algturnstile x:T}
\]
\[
\inferrule*[right=\textsc{TA-Abs}]
 {\Gamma,x:T_1\algturnstile t_2:T_2}
 {\Gamma\algturnstile\lambda x:T_1.\,t_2:T_1\to T_2}
\]
\[
\inferrule*[right=\textsc{TA-App}]
 {\Gamma\algturnstile t_1:T_1\\
  T_1=T_{11}\to T_{12}\\
  \Gamma\algturnstile t_2:T_2\\
  \algturnstile T_2\subtype T_{11}}
 {\Gamma\algturnstile t_1\,t_2:T_{12}}
\]
\[
\inferrule*[right=\textsc{TA-Rcd}]
 {\Gamma\algturnstile t_i:T_i\quad(\text{对每个 }i)}
 {\Gamma\algturnstile\{l_i=t_i\}^{i\in1..n}:
  \{l_i:T_i\}^{i\in1..n}}
\]
\[
\inferrule*[right=\textsc{TA-Proj}]
 {\Gamma\algturnstile t_1:R_1\\
  R_1=\{l_i:T_i\}^{i\in1..n}}
 {\Gamma\algturnstile t_1.l_j:T_j}
\]
\end{minipage}
\caption{算法类型关系}
\label{fig:16-3}
\end{figure}
\FloatBarrier

\textsc{TA-App} 中的等式 \(T_1=T_{11}\to T_{12}\) 只是显式记录类型检查的
执行顺序：先计算 \(t_1\) 的类型 \(T_1\)，再检查它是否为箭头类型。也可以
删掉该等式，直接把第一项前提写成
\(\Gamma\algturnstile t_1:T_{11}\to T_{12}\)。\textsc{TA-Proj} 中的
记录类型等式作用相同。应用规则的子类型前提使用算法关系；由
\cref{prop:16.1.5}，改写为声明式关系也没有区别。

\begin{exercise}[\(\star\star\), \(\nrightarrow\)]
\label{ex:16.2.3}
找出项 \(s,t\) 和类型 \(S,T\)，使
\(\algturnstile s:S\)、\(\algturnstile t:T\)、
\(s\transmulti t\)、\(T\subtype S\)，但 \(S\not\subtype T\)。这说明项在
求值后由算法规则算出的类型可能变得更小。
\end{exercise}
\taplsolutionlink{sol:translator:16.2.3}{16.2.3}

\translatornote{这里的 \(S\) 与 \(T\) 分别是两个不同项 \(s\) 与 \(t\) 的
最小类型。求值后得到的项 \(t\) 可能暴露出更多类型信息，所以它的最小类型
\(T\) 可以严格小于 \(s\) 的最小类型 \(S\)。这并不破坏声明式关系的保持性：
由 \(T\subtype S\) 和 \textsc{T-Sub}，声明式关系仍可给 \(t\) 赋予原来的类型
\(S\)。发生变化的只是算法为求值结果算出的最精确类型。}

下面正式证明算法规则与声明式规则对应。上面的推导重排只用于说明思路；把它
发展成完整证明会很冗长，直接对推导作归纳更简单。

\begin{theorem}[可靠性]
\label{thm:16.2.4}
若 \(\Gamma\algturnstile t:T\)，则 \(\Gamma\vdash t:T\)。
\end{theorem}
\begin{proof}
对算法类型推导作直接归纳。
\end{proof}

完备性的陈述稍有不同。声明式关系可为同一项赋予许多类型，算法关系至多赋予
一个类型，所以不能简单把\cref{thm:16.2.4}倒过来。可以证明：算法规则为
\(t\) 算出的唯一类型 \(S\)，是声明式规则能够赋予 \(t\) 的任意类型 \(T\)
的子类型，即 \(S\subtype T\)。再由可靠性定理，\(S\) 本身也是声明式关系
认可的类型。因此，\(S\) 位于 \(t\) 的所有可赋类型之下，正是 \(t\) 的
\termfirst{最小类型}{minimal type}。

\begin{theorem}[完备性，也称最小类型定理]
\label{thm:16.2.5}
若 \(\Gamma\vdash t:T\)，则存在 \(S\subtype T\)，使
\(\Gamma\algturnstile t:S\)。
\end{theorem}
\begin{proof}
练习（推荐，\(\star\star\)）。
\end{proof}
\taplauthorsolutionlink{sol:author:16.2.5}{16.2.5}

\begin{exercise}[\(\star\star\)]
\label{ex:16.2.6}
若删去箭头子类型规则 \textsc{S-Arrow}，但保留其余声明式子类型规则和类型
规则，系统是否仍具有最小类型性质？若有，请证明该性质；若没有，给出一个可赋型却
没有最小类型的项。
\end{exercise}
\taplauthorsolutionlink{sol:author:16.2.6}{16.2.6}

\section{并与交}
\label{sec:16.3}

带子类型化的语言在检查条件式、case 表达式等多分支构造时，还需要额外机制。
声明式条件类型规则要求两个结果分支具有相同类型：
\[
\inferrule*[right=\textsc{T-If}]
 {\Gamma\vdash t_1:\tapltype{Bool}\\
  \Gamma\vdash t_2:T\\
  \Gamma\vdash t_3:T}
 {\Gamma\vdash
  \mathsf{if}\ t_1\ \mathsf{then}\ t_2\ \mathsf{else}\ t_3:T}
\]
有包容规则时，让两个分支取得相同类型的方法可能很多。例如
\[
\mathsf{if}\ \tapltrue\
\mathsf{then}\ \{x=\tapltrue,y=\taplfalse\}\
\mathsf{else}\ \{x=\taplfalse,z=\tapltrue\}
\]
具有类型 \(\{x:\tapltype{Bool}\}\)。then 分支的最小类型为
\(\{x:\tapltype{Bool},y:\tapltype{Bool}\}\)，else 分支的最小类型为
\(\{x:\tapltype{Bool},z:\tapltype{Bool}\}\)，二者都能由包容规则提升到
\(\{x:\tapltype{Bool}\}\)。整项也能具有
\(\{x:\tapltype{Top}\}\)、\(\{\}\) 等任意公共超类型；其中最小的公共超类型
是 \(\{x:\tapltype{Bool}\}\)。

因此，计算条件式的最小类型时，应先计算两个分支的最小类型，再计算二者的
最小公共超类型。偏序论通常把它称为两种类型的
并。

\begin{definition}
\label{def:16.3.1}
若 \(S\subtype J\)、\(T\subtype J\)，并且对每个 \(U\)，从
\(S\subtype U\) 和 \(T\subtype U\) 都能推出 \(J\subtype U\)，则称
\(J\) 为 \(S,T\) 的\emph{并}，写作 \(S\sqcup T=J\)。

对偶地，若 \(M\subtype S\)、\(M\subtype T\)，并且对每个 \(L\)，从
\(L\subtype S\) 和 \(L\subtype T\) 都能推出 \(L\subtype M\)，则称
\(M\) 为 \(S,T\) 的交，写作 \(S\sqcap T=M\)。
\end{definition}

一种子类型关系不一定为每对类型都提供并或交。若每对 \(S,T\) 都有某个并，
就说该关系\emph{具有并}；若每对类型都有某个交，就说它\emph{具有交}。

本节的子类型关系\footnote{即\cref{fig:15-1,fig:15-3}定义的关系，再加入
\(\tapltype{Bool}\)。没有为 \(\tapltype{Bool}\) 增加声明式子类型规则，
所以它唯一的严格超类型是 \(\tapltype{Top}\)。}具有并，却不具有交。例如，\(\{\}\) 与
\(\tapltype{Top}\to\tapltype{Top}\) 连公共子类型都没有，自然也没有最大
公共子类型。不过，若只要求那些拥有共同下界的类型对存在交，这一较弱的性质仍然成立。
若存在 \(L\)，使
\(L\subtype S\) 且 \(L\subtype T\)，就称 \(S,T\)
\termfirst{有共同下界}{bounded below}。若每对有共同下界的类型都有交，
就说子类型关系具有\termfirst{有界交}{bounded meets}。

并和交不一定在语法上唯一。例如，
\(\{x:\tapltype{Top},y:\tapltype{Top}\}\) 与
\(\{y:\tapltype{Top},x:\tapltype{Top}\}\) 都是
\(\{x:\tapltype{Top},y:\tapltype{Top},z:\tapltype{Top}\}\) 和
\(\{x:\tapltype{Top},y:\tapltype{Top},w:\tapltype{Top}\}\) 的并。
不过，对于固定的一对类型，若 \(J_1\) 和 \(J_2\) 都是它们的并，则
\(J_1\subtype J_2\) 且 \(J_2\subtype J_1\)；交也具有同样的性质。

\begin{proposition}[并与有界交的存在性]
\label{prop:16.3.2}
\begin{enumerate}
\item 对每对类型 \(S,T\)，都存在 \(J\)，使 \(S\sqcup T=J\)。
\item 对每对具有公共子类型的 \(S,T\)，都存在 \(M\)，使
  \(S\sqcap T=M\)。
\end{enumerate}
\end{proposition}
\begin{proof}
练习（推荐，\(\star\star\star\)）。
\end{proof}
\taplauthorsolutionlink{sol:author:16.3.2}{16.3.2}

利用并运算，可以给条件式写出算法类型规则：
\[
\inferrule*[right=\textsc{TA-If}]
 {\Gamma\algturnstile t_1:T_1\\T_1=\tapltype{Bool}\\
  \Gamma\algturnstile t_2:T_2\\
  \Gamma\algturnstile t_3:T_3\\
  T_2\sqcup T_3=T}
 {\Gamma\algturnstile
  \mathsf{if}\ t_1\ \mathsf{then}\ t_2\ \mathsf{else}\ t_3:T}
\]

\begin{exercise}[\(\star\star\)]
\label{ex:16.3.3}
\(\mathsf{if}\ \tapltrue\ \mathsf{then}\ \taplfalse\
\mathsf{else}\ \{\}\) 的最小类型是什么？这是我们想要的结果吗？
\end{exercise}
\taplauthorsolutionlink{sol:author:16.3.3}{16.3.3}

\begin{exercise}[\(\star\star\star\)]
\label{ex:16.3.4}
把计算并与交的算法扩展到\taplsecref{15.5}{sec:15.5}所述那类带引用的命令式语言中，是否
容易？若再把不变的 \(\tapltype{Ref}\) 分解为协变的
\(\tapltype{Source}\) 与逆变的 \(\tapltype{Sink}\)，情况如何？
\end{exercise}
\taplauthorsolutionlink{sol:author:16.3.4}{16.3.4}

\section{算法类型关系与 \(\tapltype{Bot}\) 类型}
\label{sec:16.4}

若子类型关系加入最小类型 \(\tapltype{Bot}\)，子类型算法和类型算法也要稍作
扩展。算法子类型关系直接加入以下规则：
\[
\inferrule*[right=\textsc{SA-Bot}]
 { }{\algturnstile\tapltype{Bot}\subtype T}
\]
算法类型关系还要加入两条稍微特殊的规则：
\[
\inferrule*[right=\textsc{TA-AppBot}]
 {\Gamma\algturnstile t_1:T_1\\T_1=\tapltype{Bot}\\
  \Gamma\algturnstile t_2:T_2}
 {\Gamma\algturnstile t_1\,t_2:\tapltype{Bot}}
\]
\[
\inferrule*[right=\textsc{TA-ProjBot}]
 {\Gamma\algturnstile t_1:R_1\\R_1=\tapltype{Bot}}
 {\Gamma\algturnstile t_1.l_i:\tapltype{Bot}}
\]
在声明式系统中，包容规则可以把 \(\tapltype{Bot}\) 提升为任意函数类型，
所以一个 \(\tapltype{Bot}\) 型项可以应用于任意类型的实参，结果又能被看作
任意类型。投影的情形相同。上面两条算法规则把这种行为直接写了出来。

\begin{exercise}[\(\star\)]
\label{ex:16.4.1}
若语言还包含条件式，是否需要再为 \(\mathsf{if}\) 加一条算法类型规则？
\end{exercise}
\taplauthorsolutionlink{sol:author:16.4.1}{16.4.1}

在当前语言中，为 \(\tapltype{Bot}\) 增加这些支持并不复杂。不过，
\taplsecref{28.8}{sec:28.8}会看到，\(\tapltype{Bot}\) 与有界量化结合时
会产生更严重的困难。
